The language includes a small set of primitive types: \texttt{Int}, \texttt{Bool}, \texttt{String}, and the unit type \texttt{Unit}. These are simple by themselves, but can be combined using tuples and function types to create more powerful abstractions.

\begin{verbatim}
t ::= Int | Bool | String | Unit   (* primitive types  *)
    | TypeName                     (* named types *)
    | t1 * t2                      (* tuple types      *)
    | t1 -> t2                     (* function types   *)
    | Later t                      (* later type       *)
    | Delayed t                    (* delay type       *)
    | Signal t                     (* signal type      *)
    | t1 (t2, ..., tn)             (* type application *) 
\end{verbatim}

\subsubsection{Modal Types: Later and Delayed}

% reference back to \ref{section:frp:rizzo}

The modal types \texttt{Later} and \texttt{Delayed} are what distinguish Rizzo from ordinary functional languages. As introduced in Section~\ref{section:frp:rizzo}, these types encode temporal properties of values.
\citeauthor*{rizzo_modal_2025} introduces a set of constructors for these modal types, for which the constructors for \texttt{Later} values are:

\begin{itemize}
    \item $\mathsf{never} : \texttt{Later}\,A$ is the delayed computation that never produces any value.
    % (* a computation that never produces a value *)
    \item $\mathsf{wait}\;c : \texttt{Later}\,A$ waits for an input on channel $c : \texttt{Chan}\,A$ to produce a value.
    % (* waits for input on channel \texttt{c : Chan a} *)
    \item $\mathsf{watch}\;s : \texttt{Later}\,A$ watches a partial signal $s : \texttt{Signal}\,(\texttt{Option}\,a)$ and will produce a value when the head of $s$ is updated to a value of the form $\mathsf{Some}(x)$ for some $x: A$. 
    % \texttt{watch s~:~Later a} (* watches a partial signal $\texttt{Signal }(\texttt{Option }a)$ *)
    \item $\mathsf{tail}\;s : \texttt{Later}\,(\texttt{Signal}\,A)$ builds a delayed computation to access the future value of a signal $s : \texttt{Signal}\,A$.
    % \texttt{tail s : Later (Signal a)}   (* gets the tail of signal *)
    \item $\mathsf{laterapp} \; f \; x : \texttt{Later}\,B$ is the applicative action on \texttt{Later} values. Builds a delayed computation to apply $f : \texttt{Delayed}\,(A \to B)$ to $x : \texttt{Later}\,A$ in the future time step when $x$ produces its value.
    % \texttt{laterapp f x}  (* applies \texttt{f~:~Delayed ($a \to b$)} to \texttt{x~:~Later $a$} whenever \texttt{x} becomes available *)
    \item $\mathsf{sync}\; x \; y : \texttt{Later}\,(\texttt{Sync}\,(A,B))$ is the primitive for handling the asynchronous nature of Rizzo. This constructor builds a delayed computation that will produce a value when either $x : \texttt{Later}\,A$ or $y : \texttt{Later}\,B$, or both produce a value. This either-or behaviour is captured by the $\texttt{Sync}\,(A,B)$ type with its constructors $\mathsf{Left}(x')$, $\mathsf{Right}(y')$, and $\mathsf{Both}(x',y')$.
    % \texttt{sync l1 l2 : Later (Sync a b)}   (* synchronizes two later's, producing a value when either or both tick *)
\end{itemize}

and the constructors for the \texttt{Delayed} type are:
\begin{itemize}
    \item $\mathsf{delay}\;e : \texttt{Delayed}\,A$ delays the expression $e$ to a future time step.
    % \texttt{delay e : Delayed a}   (* creates a delayed computation from expression \texttt{e} *)
    \item $\mathsf{ostar} \; f \; x : \texttt{Delayed}\,B$ builds a delayed computation that applies $f: \texttt{Delayed}\,(A \to B)$ to $x : \texttt{Delayed}\,A$ in any future time step. 
    % \texttt{ostar f x}  (* applies \texttt{f~:~Delayed ($a \to b$)} to a delayed value \texttt{x~:~Delayed a} \todo{when it is advanced in a future time step} *)
\end{itemize}


% \paragraph{The Delayed Type.} A value of type \texttt{Delayed a} represents a computation that will be available at \emph{all} future time steps. This universally quantified nature makes \texttt{Delayed} values safe to use across time boundaries. The main constructor is \texttt{delay e : Delayed a}, which creates a delayed computation from an expression \texttt{e}. The expression \texttt{e} is not evaluated immediately, but is instead, suspended until the delayed value is applied to a later value.

% \subsubsection{The Sync Type}

% In \citetitle{rizzo_modal_2025} \cite{rizzo_modal_2025} the \texttt{sum} type is needed to represent the result of synchronizing two \texttt{Later} computations with the \texttt{sync} signal combinator. To simplify the implementation, we have chosen to use a custom \texttt{Sync} type instead, which is a three-way sum indicating which computations produced values when synchronizing two \texttt{Later} computations with \texttt{sync}:

% \begin{verbatim}
% type Sync a b = Left a | Right b | Both (a * b)
% \end{verbatim}

% In practice, the \texttt{sync} combinator is used to combine two \texttt{Later} computations, and the resulting \texttt{Sync} value indicates whether the first, second, or both signals have ticked in the current time step. This allows combinators that use \texttt{sync} to react differently depending on which inputs updated. For example, the \texttt{zip} combinator uses \texttt{sync} to combine two signals:

% \begin{verbatim}
% fun zip xs ys : Signal 'a -> Signal 'b -> Signal ('a * 'b) =
%     let cont = fun mysync ->
%         match mysync with
%         | Left (xs_)      -> zip xs_ ys
%         | Right (ys_)     -> zip xs ys_
%         | Both (xs_, ys_) -> zip xs_ ys_
%     in
%     (head xs, head ys) :: (cont |> sync (tail xs) (tail ys))
% \end{verbatim}